%! Representing a Quantitative Variable with Graphs — Unit 1, Lesson 1.3 — Student Packet Cover Sheet
\documentclass[12pt]{article}
\usepackage{apstats-article}
\usepackage{apstats-boxes}
\usepackage{ltablex}
\keepXColumns

\begin{document}

% Full-bleed navy banner. \coverbanner MEASURES the title block and sizes the
% band to it, so a lesson title long enough to wrap grows the band rather than
% running out the bottom of it.
\coverbanner{Unit 1}{Lesson 1.3 \quad Representing a Quantitative Variable with Graphs}

\namedateperiod
\vspace{0.18in}

\boxguard[14]
\begin{learningtargetbox}
\small
% GRADUAL RELEASE: the vocabulary is taught in the guided notes, so the targets name it
% plainly. The AP Learning Objectives (UNC-1.F, UNC-1.G) stay in the teacher-facing plan.
\begin{enumerate}[leftmargin=1.5em, itemsep=5pt]
  \item \ldots{} classify a quantitative variable as \textbf{discrete} or \textbf{continuous},
        and say why rounding the recorded values does not change the answer.
  \item \ldots{} read a \textbf{dotplot}, a \textbf{stem plot}, and a \textbf{histogram}, and
        say exactly what each one keeps and what it throws away.
  \item \ldots{} explain why the bars of a histogram touch when the bars of last lesson's bar
        graph did not.
  \item \ldots{} explain how changing the \textbf{interval width} changes what a histogram
        shows --- without changing a single data value.
\end{enumerate}
\end{learningtargetbox}

\vspace{0.14in}

\boxguard[14]
\begin{tocbox}
\small
\renewcommand{\arraystretch}{1.5}
\begin{tabularx}{\linewidth}{@{} c l X r @{}}
\rowcolor{sky}
\textbf{\#} & \textbf{Component} & \textbf{Description} & \textbf{Score} \\
\midrule
1 & Warm-Up & Shares, a tally, and four things that might come out to 7.5 & \blank{1.2cm} \\
2 & Guided Notes \& Practice & Two kinds of quantitative variable, three ways to draw one, and
                 the same 24 finishing times telling two different stories --- I show you, we do
                 one together, you do the rest on your own & \blank{1.2cm} \\
% AP Practice is assigned but NOT scored: its score cell prints NA, never a \blank{}.
% Row 4 is the lesson's GRADED practice. Paper homework is the DEFAULT for this course; on the
% occasions DeltaMath is assigned instead, that replaces row 4 and its score cell is struck.
3 & AP Practice & A rainfall record: AP-style multiple choice and one free-response set ---
                 practice, not a grade & \textbf{NA} \\
4 & Homework & Discrete vs.\ continuous, two histograms of one data set, a stem plot, and
                 one multiple-choice item --- \textbf{scored, due next class} & \blank{1.2cm} \\
\midrule
  & \multicolumn{2}{r}{\textbf{Total}} & \blank{1.2cm} \\
\end{tabularx}
\end{tocbox}

\vspace{0.14in}

\boxguard[8]
\begin{remindbox}
\small
Today runs in two moves inside the guided notes --- \textbf{I show you}, then \textbf{we do one
together} --- and a third at the end: \textbf{you start the homework here, on your own}, before
you leave. The notes give you the vocabulary, the three displays, and what happens when the
interval width changes; the homework is where you run it by yourself. \textbf{Two of the
pictures in this packet are both correct and disagree with each other} --- when that happens, do
not pick a favourite. Work out what was changed.
\end{remindbox}

\end{document}
